\chapter{Requirements Analysis}
\label{ch:analysis}

This chapter focuses on requirement identification, elicitation, and analysis. The users are categorized into classes. For requirement identification and elicitation, open interviews, ethnography, and use case analysis are used. The identified requirements are categorized into functional and non-functional requirements along with their use case descriptions and diagrams.

\section{User Classes and Characterization}
\label{sec:user-classes}

\begin{table}[H]
\centering
\caption{User Classes and Characteristics}
\label{tab:user-classes}
\begin{tabularx}{\textwidth}{|l|X|}
\hline
\textbf{User Class} & \textbf{Characterization} \\ \hline
Patient & The patient logs into the application to manage medications, health measurements, and appointment reminders. \\ \hline
Caregiver & The caregiver is an authorized person who can log in to manage reminders and health tasks on behalf of the patient when the patient is unavailable. \\ \hline
Doctor & The doctor logs into the application to view patient profiles after appointments and analyze their health data. The doctor gives advice on medicine, measurement, activity plan, and caregiver support. \\ \hline
Pharmacist & The pharmacist receives the request through the Pharmacy Dashboard and provides the medicine and delivery cost. The pharmacist checks and accepts the delivery address to send the medicine. \\ \hline
Admin & The admin can view registered users. Verify the doctor, pharmacist, and nurse license. Users can submit feedback which can be viewed by the admin. \\ \hline
\end{tabularx}
\end{table}

\section{Requirement Identifying Techniques}
\label{sec:req-techniques}

\subsection{Interview}
Open interviews were conducted with individuals from different age groups to understand the difficulty they face in managing medicines, appointments, and daily health activities. The application helps to identify common problems such as forgetting doses and missing appointments, which helps in defining the system requirements.

\subsection{Ethnography or Observation}
Software systems are not used alone---they work in social and organizational environments. User needs are often affected by their daily routines and surroundings. Ethnography is an observational technique that can be used to understand operational processes and help derive requirements for software to support these processes. It was noticed that many people rely on memory or simple alarms and often forget important health tasks. This observation helped in identifying the need for a structured and organized healthcare reminder system.

\section{Use Case Model}
\label{sec:use-case-model}

% Include your use case diagram images here
% \begin{figure}[H]
% \centering
% \includegraphics[width=\textwidth]{use_case_diagram_1.png}
% \caption{Use Case Diagram 1}
% \label{fig:use-case-1}
% \end{figure}

\textit{(Include Use Case Diagram images here.)}

\section{Use Case Descriptions}
\label{sec:use-case-desc}

% Use Case Descriptions - All 24 Use Cases

\begin{table}[H]
\centering
\caption{Use Case 1: Sign Up}
\begin{tabularx}{\textwidth}{|l|X|}
\hline
\textbf{Use Case ID} & UC-1 \\ \hline
\textbf{Use Case Name} & Sign Up \\ \hline
\textbf{Actors} & Patient/Caregiver/Doctor/Pharmacist \\ \hline
\textbf{Description} & User creates a new account by filling the sign-up form. \\ \hline
\textbf{Trigger} & User opens the app and selects Sign Up. \\ \hline
\textbf{Preconditions} & Internet connection is available. \\ \hline
\textbf{Postconditions} & User's account is created and stored in the database. \\ \hline
\textbf{Normal Flow} & 1. User opens app and clicks Sign Up. 2. Enters name, email, DOB, gender, password. 3. App validates and creates account. 4. Verification email sent. 5. User redirected to login. \\ \hline
\textbf{Alternative Flows} & If email already exists, error is shown. \\ \hline
\textbf{Business Rules} & Email must be unique and valid. Password must meet security requirements. \\ \hline
\textbf{Assumptions} & Internet is available. \\ \hline
\end{tabularx}
\end{table}

\begin{table}[H]
\centering
\caption{Use Case 2: Log In}
\begin{tabularx}{\textwidth}{|l|X|}
\hline
\textbf{Use Case ID} & UC-2 \\ \hline
\textbf{Use Case Name} & Log In \\ \hline
\textbf{Actors} & Patient/Caregiver/Doctor/Pharmacist \\ \hline
\textbf{Description} & Allows user to access the system by entering email and password. \\ \hline
\textbf{Trigger} & User opens the app and selects Log In. \\ \hline
\textbf{Preconditions} & User must have a registered account. Internet connection is available. \\ \hline
\textbf{Postconditions} & User is authenticated and redirected to the home screen. \\ \hline
\textbf{Normal Flow} & 1. User opens login screen. 2. Enters email and password. 3. System validates credentials. 4. User is redirected to Home screen. \\ \hline
\textbf{Alternative Flows} & If credentials are wrong, show error message. If email is not verified, ask user to verify. \\ \hline
\textbf{Business Rules} & Only registered users can log in. \\ \hline
\textbf{Assumptions} & Internet is available. \\ \hline
\end{tabularx}
\end{table}

\begin{table}[H]
\centering
\caption{Use Case 3: Forgot Password}
\begin{tabularx}{\textwidth}{|l|X|}
\hline
\textbf{Use Case ID} & UC-3 \\ \hline
\textbf{Use Case Name} & Forgot Password \\ \hline
\textbf{Actors} & Patient/Caregiver/Doctor/Pharmacist \\ \hline
\textbf{Description} & Allows user to recover password by entering registered email. \\ \hline
\textbf{Trigger} & User clicks ``Forgot Password'' on login screen. \\ \hline
\textbf{Preconditions} & User must have a registered account. \\ \hline
\textbf{Postconditions} & A password reset link is sent to the user's email. \\ \hline
\textbf{Normal Flow} & 1. User selects Forgot Password. 2. Enters registered email. 3. System sends reset link. 4. User resets password via email. \\ \hline
\textbf{Alternative Flows} & If email is not registered, system shows error. \\ \hline
\textbf{Business Rules} & Email must be registered. Reset link expires after limited time. \\ \hline
\textbf{Assumptions} & Internet and email access available. \\ \hline
\end{tabularx}
\end{table}

\begin{table}[H]
\centering
\caption{Use Case 4: Add Medication}
\begin{tabularx}{\textwidth}{|l|X|}
\hline
\textbf{Use Case ID} & UC-4 \\ \hline
\textbf{Use Case Name} & Add Medication \\ \hline
\textbf{Actors} & Patient/Caregiver \\ \hline
\textbf{Description} & User adds medication details to set up a reminder. \\ \hline
\textbf{Trigger} & User navigates to the Medication section and taps Add. \\ \hline
\textbf{Preconditions} & User must be logged in. \\ \hline
\textbf{Postconditions} & Medication is saved and a reminder is scheduled. \\ \hline
\textbf{Normal Flow} & 1. User opens Medication screen. 2. Enters medicine name, dosage, frequency, start/end date, reminder time. 3. System validates and saves. 4. Reminder notification is scheduled. \\ \hline
\textbf{Alternative Flows} & If required fields are empty, show error message. \\ \hline
\textbf{Business Rules} & Medicine name cannot be empty. Start date cannot be later than end date. \\ \hline
\textbf{Assumptions} & Internet is available. \\ \hline
\end{tabularx}
\end{table}

\begin{table}[H]
\centering
\caption{Use Case 5: Set Measurement Reminder}
\begin{tabularx}{\textwidth}{|l|X|}
\hline
\textbf{Use Case ID} & UC-5 \\ \hline
\textbf{Use Case Name} & Set Measurement Reminder \\ \hline
\textbf{Actors} & Patient/Caregiver \\ \hline
\textbf{Description} & User sets a reminder for health measurement (BP, sugar, weight). \\ \hline
\textbf{Trigger} & User navigates to Measurement section and selects Add Reminder. \\ \hline
\textbf{Preconditions} & User must be logged in. \\ \hline
\textbf{Postconditions} & Measurement reminder is saved and notification scheduled. \\ \hline
\textbf{Normal Flow} & 1. User selects measurement type. 2. Sets reminder time. 3. System validates and saves. 4. Notification is scheduled. \\ \hline
\textbf{Alternative Flows} & If type or time is not selected, show error. \\ \hline
\textbf{Business Rules} & Measurement type must be selected. Reminder time must be valid. \\ \hline
\textbf{Assumptions} & Internet is available. \\ \hline
\end{tabularx}
\end{table}

\begin{table}[H]
\centering
\caption{Use Case 6: Set Activity Reminder}
\begin{tabularx}{\textwidth}{|l|X|}
\hline
\textbf{Use Case ID} & UC-6 \\ \hline
\textbf{Use Case Name} & Set Activity Reminder \\ \hline
\textbf{Actors} & Patient/Caregiver \\ \hline
\textbf{Description} & User sets reminders for daily activities (walking, water, meals, exercise, yoga, cycling). \\ \hline
\textbf{Trigger} & User navigates to Activity section. \\ \hline
\textbf{Preconditions} & User must be logged in. \\ \hline
\textbf{Postconditions} & Activity reminder is saved and notification scheduled. \\ \hline
\textbf{Normal Flow} & 1. User selects activity type. 2. Sets reminder time and frequency. 3. System validates and saves. 4. Notification is scheduled. \\ \hline
\textbf{Alternative Flows} & If fields are empty, show error. \\ \hline
\textbf{Business Rules} & Activity type must be selected. \\ \hline
\textbf{Assumptions} & Internet is available. \\ \hline
\end{tabularx}
\end{table}

% NOTE: Use Cases 7-24 follow the same tabular format.
% Due to space, they are summarized here. In your final document,
% replicate the table structure above for each use case:
% UC-7: Send Request to Doctor
% UC-8: Send Request to Caregiver
% UC-9: Doctor Module (View Patient Data, Give Recommendations)
% UC-10: Doctor Recommendation
% UC-11: In-App Chat (Doctor-Patient)
% UC-12: In-App Chat (Patient-Caregiver)
% UC-13: Doctor Payment
% UC-14: Caregiver Module
% UC-15: Set Appointment Reminder
% UC-16: Refill Reminder
% UC-17: Measurement Tracker
% UC-18: Pharmacy Medicine Delivery
% UC-19: Chatbot Assistance
% UC-20: Track Intake
% UC-21: View Profile
% UC-22: Admin Module
% UC-23: Feedback
% UC-24: Emergency Contact

\begin{table}[H]
\centering
\caption{Use Case 7: Send Request to Doctor}
\begin{tabularx}{\textwidth}{|l|X|}
\hline
\textbf{Use Case ID} & UC-7 \\ \hline
\textbf{Use Case Name} & Send Request to Doctor \\ \hline
\textbf{Actors} & Patient \\ \hline
\textbf{Description} & The patient sends a connection request to a doctor by entering doctor details and symptoms. \\ \hline
\textbf{Trigger} & Patient selects Doctor tab in Medical Professional section. \\ \hline
\textbf{Preconditions} & Patient must be logged in. Doctor must be registered. \\ \hline
\textbf{Postconditions} & Request is sent and status is set to Pending. \\ \hline
\textbf{Normal Flow} & 1. Patient opens Doctor section. 2. Enters doctor name, email, speciality, symptoms. 3. Sends request. 4. System saves with Pending status. \\ \hline
\textbf{Alternative Flows} & If doctor email is not found, error shown. \\ \hline
\textbf{Business Rules} & All required fields must be filled. \\ \hline
\textbf{Assumptions} & Internet is available. \\ \hline
\end{tabularx}
\end{table}

\begin{table}[H]
\centering
\caption{Use Case 8: Send Request to Caregiver}
\begin{tabularx}{\textwidth}{|l|X|}
\hline
\textbf{Use Case ID} & UC-8 \\ \hline
\textbf{Use Case Name} & Send Request to Caregiver \\ \hline
\textbf{Actors} & Patient \\ \hline
\textbf{Description} & The patient sends a request to a caregiver by entering their name and email. \\ \hline
\textbf{Trigger} & Patient selects Caregiver tab in Medical Professional section. \\ \hline
\textbf{Preconditions} & Patient must be logged in. Caregiver must be registered. \\ \hline
\textbf{Postconditions} & Request sent and status set to Pending. \\ \hline
\textbf{Normal Flow} & 1. Patient opens Caregiver section. 2. Enters caregiver name and email. 3. Sends request. 4. System saves with Pending status. \\ \hline
\textbf{Alternative Flows} & If caregiver email not found, error shown. \\ \hline
\textbf{Business Rules} & All required fields must be filled. \\ \hline
\textbf{Assumptions} & Internet is available. \\ \hline
\end{tabularx}
\end{table}

% Continue with remaining use cases in the same format...
% UC-9 through UC-24 should be added following the same table pattern.


\section{Functional Requirements}
\label{sec:functional-req}

This section describes the functional requirements of the system expressed in natural language style. This section is organized by features as a system feature name and specific functional requirements associated with this feature.

% Functional Requirements - FR-1 through FR-24

\begin{table}[H]
\centering
\caption{Functional Requirement 1: User Profile Management}
\begin{tabularx}{\textwidth}{|l|X|}
\hline
\textbf{Identifier} & FR-1 \\ \hline
\textbf{Title} & User Profile Management \\ \hline
\textbf{Requirement} & The user shall be able to create and manage their profile by entering personal information such as full name, email, Date of Birth, gender, and password. When a user selects their role, it will also enter a license number or speciality. The user shall be able to login by entering email and password. \\ \hline
\textbf{Source} & User Requirement, System Requirement \\ \hline
\textbf{Rationale} & It provides a secure environment and collects profile data. \\ \hline
\textbf{Business Rule} & User Date of Birth must follow a format (YYYY/MM/DD). Gender selection must be from predefined options (Male, Female, Others). \\ \hline
\textbf{Dependencies} & Internet connection, Firebase Authentication \\ \hline
\textbf{Priority} & Medium \\ \hline
\end{tabularx}
\end{table}

\begin{table}[H]
\centering
\caption{Functional Requirement 2: Forgot Password}
\begin{tabularx}{\textwidth}{|l|X|}
\hline
\textbf{Identifier} & FR-2 \\ \hline
\textbf{Title} & Forgot Password \\ \hline
\textbf{Requirement} & The end user will be able to recover the password by entering a registered email and receiving a password reset link. \\ \hline
\textbf{Source} & System Requirement \\ \hline
\textbf{Rationale} & Support work will be reduced and user will recover their password easily. \\ \hline
\textbf{Business Rule} & Password should expire after a limited time. \\ \hline
\textbf{Dependencies} & FR-1 \\ \hline
\textbf{Priority} & Medium \\ \hline
\end{tabularx}
\end{table}

\begin{table}[H]
\centering
\caption{Functional Requirement 3: Patient Health Management}
\begin{tabularx}{\textwidth}{|l|X|}
\hline
\textbf{Identifier} & FR-3 \\ \hline
\textbf{Title} & Patient Health Management \\ \hline
\textbf{Requirement} & The system allows patients to manage medications, health data, daily activities, and appointments. It supports intake updates, reminders, pharmacy, chatbot access, history, emergency contacts and feedback. \\ \hline
\textbf{Source} & User Requirement \\ \hline
\textbf{Rationale} & This feature helps patients manage their health efficiently in one place, reduces missed medications, and improves overall care and monitoring. \\ \hline
\textbf{Business Rule} & Record tasks, validate numeric data and send reminders. \\ \hline
\textbf{Dependencies} & Authentication, database, notifications and internet connection \\ \hline
\textbf{Priority} & High \\ \hline
\end{tabularx}
\end{table}

% NOTE: FR-4 through FR-24 follow the same tabular format.
% Add them following the same pattern as above.
% FR-4: Caregiver Assistance (High)
% FR-5: Pharmacist Medicine and Delivery Management (High)
% FR-6: Doctor (High)
% FR-7: In-app Chat (High)
% FR-8: Doctor's Payment (High)
% FR-9: Add Medication (High)
% FR-10: Set Measurement Reminder (High)
% FR-11: Set Activity Reminder (Medium)
% FR-12: Medical Professional (High)
% FR-13: Appointment Reminder (High)
% FR-14: Medicine Refill Reminder (High)
% FR-15: Edit Health Task (High)
% FR-16: Record Measurement Data (High)
% FR-17: Pharmacy Medicine Delivery (Medium)
% FR-18: Chatbot Assistance (Medium)
% FR-19: Receive Notification (Medium)
% FR-20: Track Intake (High)
% FR-21: View User Profile (High)
% FR-22: Admin (High)
% FR-23: Feedback (High)
% FR-24: Emergency Contact (High)


\section{Non-Functional Requirements}
\label{sec:non-functional-req}

Non-functional requirements are the constraints that specify the quality attributes of software. They are important because they can significantly impact the overall quality and success of a software system.

\subsection{Reliability}
The Smart Healthcare Reminder application should work smoothly and be available most of the time. The application should recover automatically if small errors occur and keep important data safe like reminders, intake tracking, and user profiles.

\subsection{Usability}
The application should be easy to use for all users including patients, caregivers, doctors, pharmacists, and administrators. The system should have a simple and clear interface so users should easily perform tasks like adding reminders, tracking intake, viewing profiles, and using the chatbot.

\subsection{Performance}
The application should load screens quickly under normal conditions. It should handle multiple users at the same time without slowing down. Features like saving reminders, recording measurement values, and chatbot should respond smoothly with minimal delay.

\subsection{Security}
The application should protect all user data such as personal information, health records, and reminders. Only authorized users can access the system and each user can access their own data. The system should use role-based access so that admin and users have different accessibility.

\section{External Interface Analysis}
\label{sec:external-interface}

\subsection{User Interface Requirements}
\begin{itemize}
    \item The app follows standard design principles for clarity and ease of use.
    \item It will adjust to different screen sizes and devices.
    \item Users will receive clear notifications for their actions.
\end{itemize}

\subsection{Software Interfaces}
\begin{itemize}
    \item The app will run on Android.
    \item Firebase will be used for data storage and management.
\end{itemize}

\subsection{Hardware Interfaces}
\begin{itemize}
    \item Users will access the application mainly through smartphones and tablets.
\end{itemize}

\subsection{Communication Interfaces}
\begin{itemize}
    \item The application will use standard mobile communication methods and secure protocols for data exchange.
\end{itemize}

\subsection{Error Handling and Reporting}
\begin{itemize}
    \item The application handles errors by checking user input, login details, internet connection, and notification permission. When a problem occurs, clear messages are shown to help users fix it.
\end{itemize}

\section{Summary}
\label{sec:ch2-summary}

Chapter 2 outlines the requirements of the Smart Healthcare Reminder application, detailing how users can register, set reminders for medication, measurement, activity, appointment, and refills, and track their daily intake. This chapter also outlines non-functional requirements such as reliability, usability, and security along with system interfaces used for notification, data storage, and user interaction.
